\documentclass[11pt]{article}

\usepackage[margin=1in]{geometry}
\usepackage{microtype}
\usepackage{graphicx}
\usepackage{subcaption}
\usepackage{booktabs}
\usepackage{hyperref}
\usepackage{algorithm}
\usepackage{algorithmic}

\usepackage{amsmath,amssymb,mathtools,amsthm}
\usepackage[capitalize,noabbrev]{cleveref}
\usepackage[numbers,sort&compress]{natbib}

\hypersetup{
  colorlinks=true,
  linkcolor=blue,
  citecolor=blue,
  urlcolor=blue
}

%%%%%%%%%%%%%%%%%%%%%%%%%%%%%%%%
% THEOREMS
%%%%%%%%%%%%%%%%%%%%%%%%%%%%%%%%
\theoremstyle{plain}
\newtheorem{theorem}{Theorem}[section]
\newtheorem{proposition}[theorem]{Proposition}
\newtheorem{lemma}[theorem]{Lemma}
\newtheorem{corollary}[theorem]{Corollary}
\theoremstyle{definition}
\newtheorem{definition}[theorem]{Definition}
\newtheorem{assumption}[theorem]{Assumption}
\theoremstyle{remark}
\newtheorem{remark}[theorem]{Remark}

\title{Where to Search: Global Scheduling of Noise Trajectory Search in Diffusion Models}
\author{Anonymous Author(s)}
\date{}

\begin{document}

\maketitle

\begin{abstract}
Test-time compute can substantially improve diffusion model generation via noise search and verifier-guided sampling, but existing methods largely allocate search compute uniformly or with fixed heuristics along timesteps.
We study \emph{where to search} under a strict budget of function evaluations (NFE), and formalize inference-time noise refinement as a two-level problem: a \emph{local} operator that searches for better noise at a chosen timestep, and a \emph{global} scheduler that allocates per-timestep search iterations across the trajectory.
Building on this view, we propose a scheduling algorithm that combines offline profiling with online control: offline analysis provides a coarse per-step budget allocation, and an online controller performs windowed early stopping based on observed score gains and score variance, while enforcing the exact NFE budget.
Experiments on Stable Diffusion and EDM show consistent improvements and faster quality scaling: our method achieves higher scores with 20--50\% fewer NFEs at matched quality levels.
\end{abstract}

\section{Introduction}
\label{sec:intro}

Diffusion models have become a dominant paradigm for high-quality image and
multimodal generation, powering modern systems such as Stable Diffusion and
large-scale class-conditional generators.
While most progress has focused on improving training objectives and model
architectures, an emerging line of work shows that \emph{inference-time
scaling}---allocating more computation at test time without retraining---can
significantly improve sample quality under a fixed model.
Recent advances demonstrate that test-time techniques such as noise search,
trajectory refinement, and verifier-guided sampling can yield substantial gains
over plain diffusion sampling, especially when additional compute is available.

Importantly, our framework applies broadly to both stochastic and deterministic
samplers. While some diffusion models---particularly flow-based methods and
ODE solvers---use deterministic transitions after the initial noise,
inference-time search \emph{requires} stochasticity to explore different
generation paths. For deterministic samplers, this is achieved by intentionally
injecting noise at selected timesteps, enabling the search to explore diverse
trajectories that would otherwise be inaccessible.
Recent work has formalized this ODE-to-SDE conversion, showing that any
deterministic flow can be transformed into a family of SDEs with the same
marginal distributions~\citep{dockhorn2024stochastic,kim2025rbf}.
This deliberate introduction of stochasticity for exploration has been adopted
in prior work on reinforcement learning for diffusion
models~\citep{black2023training,fan2024reinforcement} and is conceptually
similar to exploration bonuses in bandit algorithms.

A key limitation of existing inference-time methods, however, is that they
largely treat the denoising trajectory as a homogeneous sequence of steps.
Most approaches either (i) allocate computation uniformly across timesteps, or
(ii) hard-code search to specific parts of the trajectory (e.g., only the
initial noise or a fixed middle window).
This overlooks a fundamental empirical fact: \emph{different timesteps exhibit
dramatically different sensitivities to noise perturbations}.
In practice, refining noise at some steps yields large downstream improvements
in sample quality, while spending the same computation at other steps has
negligible effect.
As a result, naive per-step allocation wastes a significant fraction of the
test-time compute budget.

In this work, we study \emph{where to search} along the diffusion trajectory.
We formalize inference-time noise refinement as a \emph{noise trajectory search}
problem under a fixed number of function evaluations (NFE), and ask how to
\emph{globally schedule} search computation across timesteps.
Our central observation is that effective inference-time scaling requires
reasoning at two distinct levels: \emph{local} decisions about how to search for
better noise at a given timestep, and \emph{global} decisions about how to
distribute computation across the entire denoising process.
While prior work has explored increasingly sophisticated local search operators,
the problem of global per-step compute allocation has only recently begun
to receive attention, with concurrent work proposing either uniform allocation
with implicit adaptation~\citep{ramesh2025nts} or purely reactive online
reallocation~\citep{kim2025rbf}.

To bridge this gap, we propose a unified two-level framework for noise trajectory
search in diffusion models.
At the local level, we abstract a broad family of existing noise refinement
methods as a generic local search operator applied at a single timestep.
At the global level, we introduce a scheduler that decides \emph{when} and
\emph{how much} local search to perform at each timestep under a total NFE
budget.
Within this framework, we show that standard diffusion sampling, best-of-$N$
selection, $\epsilon$-greedy noise search, and particle- or tree-based methods
can all be interpreted as special cases with fixed or heuristic scheduling
policies.

Building on this formulation, we propose a new global scheduling algorithm that
combines \emph{offline profiling} with \emph{online control}.
Offline, we analyze verifier sensitivity across timesteps to obtain a coarse
allocation of compute tailored to the model and dataset.
Online, we dynamically adjust per-step search using feedback from each instance,
such as score gains and variance, allowing the sampler to stop early on
unproductive steps and reallocate computation where it is most effective.
Crucially, this controller operates without access to future timesteps and
respects a strict NFE budget.

We evaluate our approach on Stable Diffusion, EDM, and ImageNet benchmarks under
fixed compute budgets.
Across a wide range of NFE settings and evaluation metrics, our global
scheduling strategy consistently outperforms uniform allocation and naive
online heuristics.
These results demonstrate that \emph{global scheduling with per-step awareness}
is a critical yet previously underexplored dimension of inference-time scaling
for diffusion models.

\paragraph{Contributions.}
Our main contributions are:
(i) a unified two-level view of noise trajectory search that separates local
noise refinement from global timestep scheduling;
(ii) a novel global scheduling algorithm that integrates \emph{offline timestep
profiling} with \emph{online feedback control}, combining learned difficulty
patterns per step with adaptive early stopping per instance, unlike purely
uniform~\citep{ramesh2025nts} or purely reactive~\citep{kim2025rbf} approaches;
and (iii) extensive experiments showing consistent gains over existing
inference-time methods under identical NFE budgets.


%=================================
\section{Related Work}
\label{sec:related}

\paragraph{Inference-time scaling for generative models.}
Inference-time scaling has recently emerged as a powerful alternative to
training-time improvements, enabling quality gains by spending additional
computation at test time.
In diffusion models, this includes techniques such as classifier and
classifier-free guidance, guidance rescaling, and repeated sampling with
selection.
More recent work explores search-based methods that explicitly modify the noise
trajectory during sampling, demonstrating that inference-time computation can
be traded for higher sample quality without retraining the model.
Our work builds on this line of research but focuses on a previously overlooked
aspect: how to \emph{allocate} inference-time compute across the denoising
trajectory.

\paragraph{Noise search and trajectory refinement.}
Several methods perform local search over noise variables during diffusion
sampling.
Best-of-$N$ and random search methods sample multiple trajectories and select
the best final output according to a verifier or heuristic metric, but do not
refine noise within a trajectory.
Other approaches perform iterative noise refinement or greedy search at each
timestep, often guided by a verifier or reward function~\citep{ma2024inference}.
Most recently, \citet{ramesh2025nts} cast diffusion as an MDP and relax it to
a sequence of independent contextual bandits, proposing an $\epsilon$-greedy
algorithm that implicitly adapts exploration vs.\ exploitation across timesteps;
however, they allocate \emph{uniform} search iterations $K$ per timestep and
leave adaptive $K_t$ allocation as future work.
Concurrently, \citet{kim2025rbf} introduce Rollover Budget Forcing (RBF) for
flow models, which uses early stopping when a better particle is found and
rolls over unused budget to subsequent timesteps---a purely online mechanism
without offline profiling.
While these methods differ in how candidate noises are generated and evaluated,
they typically apply search uniformly across timesteps or rely on reactive
online adaptation without exploiting difficulty patterns that vary by step.
In contrast, our framework explicitly combines \emph{offline profiling} to
learn where search matters with \emph{online control} for per-instance
adaptation, achieving non-uniform budget allocation with strict NFE guarantees.

\paragraph{Online and adaptive inference strategies.}
Adaptive inference methods adjust computation based on feedback observed during
generation.
Examples include $\epsilon$-greedy or bandit-style noise search~\citep{ramesh2025nts},
adaptive guidance scaling, and methods that stop sampling early based on
confidence or entropy criteria.
Rollover Budget Forcing~\citep{kim2025rbf} is a recent example that reallocates
unused budget across timesteps in an online fashion.
Although these approaches introduce online decision-making, they are usually
\emph{reactive}: they start from uniform allocation and adapt only based on
real-time feedback, without leveraging prior knowledge about which timesteps
are inherently more amenable to search.
Our online controller differs in that it operates within a global budget shaped
by offline timestep profiling, explicitly balancing local gains against future
opportunities while starting from an informed, non-uniform allocation.


\paragraph{Adaptive compute allocation.}
Adaptive computation has been studied extensively in other domains, including
early exiting in neural networks, adaptive depth inference, and budgeted
optimization.
In diffusion models, recent work explores adaptive timestep skipping or variable
step sizes, but these techniques primarily modify the base sampler rather than
allocating additional search computation.
Our approach is complementary: we keep the underlying sampler fixed and focus on
how to schedule inference-time search around it.
\section{Methodology}
\label{sec:method}

\subsection{Global Modeling of Noise Trajectory Search}
\label{sec:method-global}

We consider a $T$-step diffusion sampler with discrete timesteps
$t_T > t_{T-1} > \dots > t_1$.
A pretrained diffusion model (UNet + scheduler) defines a transition
\begin{equation}
  x_{t-1} = F_\theta(x_t, t, c, \varepsilon_t),
\end{equation}
where $x_t$ is the latent at timestep $t$, $c$ is the conditioning signal
(e.g., a class label or text prompt), and $\varepsilon_t$ denotes the injected
noise (variance noise or stochastic term of the sampler).
After the final step we obtain $x_0$, which is decoded into an image and
evaluated by a verifier $v(x_0, c)$ that measures sample quality or alignment
with the condition.
We emphasize that the verifier $v$ is \emph{taken as given} and is not trained
as part of our method. In practice, verifiers can take various forms depending
on the application: perceptual quality metrics (e.g., CLIP score, aesthetic
predictors), task-specific reward functions (e.g., object detection confidence),
or simple heuristics (e.g., image brightness, compressibility). This setup
mirrors prior work on reward-guided diffusion
sampling~\citep{clark2024directly,prabhudesai2023aligning,uehara2025tutorial},
where off-the-shelf reward models guide the generation process without
additional training.

We view \emph{noise trajectory search} as modifying the collection of noise
variables $\{\varepsilon_t\}_{t=1}^T$ around a latent trajectory
$\{x_t\}_{t=1}^T$ under a fixed test-time compute budget, with the goal of
improving the verifier score.
Conceptually, it is convenient to think in terms of two levels:

\paragraph{Local search operator.}
At a single timestep $t$, given a current latent $x_t$, prompt $c$, verifier $v$
and base sampler $F_\theta$, a \emph{local} search operator
\[
  \mathcal{L}_t : (x_t, c, v, F_\theta, \text{budget at } t)
  \longrightarrow (x_{t-1}, \varepsilon_t)
\]
tries to find a better noise $\varepsilon_t$ using a small per-step budget.
To evaluate a candidate noise, the diffusion model first predicts a clean image
from the current latent:
\begin{equation}
  \hat{x}_0 = D_\theta(x_t, t, c),
  \label{eq:x0_prediction}
\end{equation}
where $D_\theta$ is the denoising network. The verifier then scores this
prediction: $s = v(\hat{x}_0, c)$. Based on the predicted $\hat{x}_0$, the
transition to the next latent is computed as
$x_{t-1} = \alpha_{t-1} \hat{x}_0 + \sigma_{t-1} \varepsilon_t$,
where $\alpha_{t-1}, \sigma_{t-1}$ are scheduler coefficients.

Concretely, a local search iteration samples $K$ candidate noises
$\{\varepsilon_t^{(k)}\}_{k=1}^K$, computes the corresponding predictions
$\{\hat{x}_0^{(k)}\}_{k=1}^K$, evaluates their scores
$\{s^{(k)} = v(\hat{x}_0^{(k)}, c)\}_{k=1}^K$, and selects the best candidate
$k^* = \arg\max_k s^{(k)}$. If $s^{(k^*)}$ improves over the current best,
the noise $\varepsilon_t^{(k^*)}$ is used; otherwise, the original noise is
retained.
Many existing zero-order or $\epsilon$-greedy noise search procedures can be
seen as specific instantiations of this operator, but here we treat
$\mathcal{L}_t$ abstractly as a local search mapping from $x_t$ to
$x_{t-1}$ given a small local budget.

\paragraph{Global scheduler over timesteps.}
While the local operator $\mathcal{L}_t$ decides \emph{how} to search at a given
step once we commit to spend budget there, it does not specify \emph{where
along the denoising trajectory} we should invest computation.
We therefore define a global scheduler
\[
  \mathcal{G} :
  \bigl(x_T, c, v, F_\theta, B, \{\mathcal{L}_t\}_{t=1}^T\bigr)
  \longrightarrow \bigl(x_0, \{x_t\}_{t=1}^T\bigr),
\]
which, starting from an initial noise $x_T \sim \mathcal{N}(0, I)$, repeatedly
applies denoising steps and local search under a total compute budget $B$ and
returns both the final sample $x_0$ and the whole latent trajectory
$\{x_t\}_{t=1}^T$.

Operationally, the scheduler $\mathcal{G}$ traverses timesteps from $T$ to $1$.
At each step $t$ it decides whether to:
(i) invoke the local operator $\mathcal{L}_t$ one or multiple times, consuming
part of the remaining budget at this timestep, or
(ii) skip local search and directly apply a single base sampler step
$x_{t-1} = F_\theta(x_t, t, c, \varepsilon_t)$.
The key challenge is that this decision must be made \emph{online}: at step $t$
we do not yet know how effective search will be at later timesteps, but any
computation used now cannot be reused later.
Our goal is to design global scheduling policies that decide, under a fixed NFE
budget, \emph{on which timesteps to search and for how long}.

This two-level view --- a local operator $\mathcal{L}_t$ and a global scheduler
$\mathcal{G}$ --- provides a common language that covers many prior
inference-time search strategies:
\begin{itemize}
  \item \textbf{Plain diffusion sampling.}
  If we let $\mathcal{L}_t$ be a single base denoising step and never allocate
  extra budget, the scheduler $\mathcal{G}$ simply applies $F_\theta$ once per
  timestep, recovering the standard sampler without search.
  \item \textbf{Best-of-$N$ / random search.}
  Taking multiple i.i.d.\ initial noises, propagating each to time $0$ without
  modifying noises, and finally selecting the best $x_0$ corresponds to a
  scheduler that invests all compute into running full trajectories and only
  performs a global selection at the end; there is effectively no local search.
  \item \textbf{Noise-trajectory and tree-style search.}
  Methods that iteratively refine a noise pivot or expand/prune partial
  trajectories at certain timesteps can be interpreted as choosing specific
  forms of $\mathcal{L}_t$ and hard-coded policies in $\mathcal{G}$ that decide
  on which timesteps to apply search and how to split budget across time.
\end{itemize}
In the remainder of this section we focus on a particular and practically
important family of methods and develop a new global scheduler within this
framework.

\subsection{Proposed Global Scheduling Algorithm}
\label{sec:method-ours}

In this work we focus on the widely used setting where, at any timestep $t$,
the sampler ultimately maintains \emph{a single latent and a single noise}
$(x_t, \varepsilon_t)$, but is allowed to perform \emph{multiple local search
iterations} inside this timestep before committing to $x_{t-1}$.
Concretely, at step $t$ we run a local noise search operator $\mathcal{L}_t$
around the current latent $x_t$, and only the best candidate noise is kept and
used to produce $x_{t-1}$ via $F_\theta$.
Under this restriction, the global scheduling problem reduces to a sequence of
\emph{per-timestep decisions}: at each $t$, given the history of search so far
and the remaining budget, should we continue local search at this timestep,
or stop and move on to $t-1$?

This raises an inherent difficulty.
When making the decision at timestep $t$, we do not know:
(1) the potential improvement that additional search at $t$ could still bring,
and (2) the potential improvement that could be obtained by saving NFE and
spending it on later timesteps.
We therefore propose a scheduling strategy that tackles these two aspects
from complementary angles.

\paragraph{Offline coarse time scheduling.}
First, for each model and dataset, we perform an \emph{offline} analysis on a
small set of prompts or images.
We run noise search with a simple local operator $\mathcal{L}_t$ and a uniform,
relatively large number of iterations at every timestep, and log the verifier
scores after each local search iteration.
From these logs we can compute, for each timestep $t$, statistics such as the
average score gain brought by local search and its variability across examples.
Aggregating these statistics over many runs reveals which parts of the
denoising trajectory are consistently more sensitive to noise refinement and
thus more ``valuable'' to search.
In our offline profiling, we observe a clear model-dependent pattern:
for Stable Diffusion (SD), most of the gains concentrate on early timesteps,
whereas for EDM the improvements are concentrated in a contiguous middle segment
of the denoising trajectory.

Based on this analysis, we partition timesteps into a small number of groups,
such as a \emph{high-value} region and a \emph{low-value} region, and assign a
larger share of the total NFE budget to the high-value region.
Within each region, we derive a coarse per-step budget $\{K_t\}$ (e.g., a
higher average number of local iterations in high-value steps and a lower
average in low-value steps).
This yields a time schedule tailored to the model and dataset, fixed across
prompts, that captures \emph{where search tends to matter the most}.

\paragraph{Online fine-grained scheduling.}
Second, at test time, we add an \emph{online} controller that refines how the
pre-assigned budget $\{K_t\}$ is actually used, on a per-prompt basis.
For each timestep $t$, we run local search up to $K_t$ iterations, but allow the
controller to stop earlier or, within a small window, adaptively modulate how
aggressively to spend the budget.

The controller uses two simple statistics computed from the local search
feedback at timestep $t$:
(i) the \emph{score gain} of local search, and
(ii) the \emph{score variance} across candidates within this step.
Let $s_t^{(j)}$ be the best verifier score at step $t$ after $j$ local search
iterations, and let $\{s_{t,\text{cand}}^{(j,i)}\}_i$ denote the scores of all
candidate noises evaluated in iteration $j$.
We define the incremental gain
\[
  g_t^{(j)} = s_t^{(j)} - s_t^{(j-1)}
\]
and the per-iteration variance
\[
  \operatorname{Var}_t^{(j)} = \operatorname{Var}_i\bigl(s_{t,\text{cand}}^{(j,i)}\bigr).
\]
Intuitively, $g_t^{(j)}$ measures whether we are still finding significantly
better noises at step $t$, while $\operatorname{Var}_t^{(j)}$ reflects how
sensitive the verifier is to noise perturbations at this timestep.

Our online policy is based on the following intuition:
if the recent gains $g_t^{(j)}$ remain positive and the score variance
$\operatorname{Var}_t^{(j)}$ is large, then this timestep is both important
and not yet saturated, so we continue local search until we reach the allocated
budget $K_t$ (or a small upper margin).
Conversely, if gains have diminished (e.g., recent $g_t^{(j)}$ are close to
zero or negative) and the variance across candidates is small, then local
search at this timestep is likely exhausted and further NFE is better reserved
for subsequent timesteps; in this case we stop early and move on to $t-1$.
These decisions are made online based solely on statistics observed up to step
$t$, without access to future timesteps. To ensure the realized compute exactly matches the prescribed total NFE budget,
we reconcile any over-/under-spending caused by adaptive stopping/slack:
any extra NFE consumed in the high-value region is compensated by allocating
additional iterations to the first few low-value steps, while any NFE deficit
is corrected by trimming iterations from the last low-value steps.

\paragraph{Overall procedure.}
Putting the pieces together, a single sampling run proceeds as follows.
We first draw an initial noise $x_T \sim \mathcal{N}(0, I)$ and obtain the
coarse per-step budget $\{K_t\}$ from the offline time schedule.
Then, for $t = T, \dots, 1$, we run the local noise search operator
$\mathcal{L}_t$ around the current latent $x_t$, while the online controller
monitors the sequence of gains $\{g_t^{(j)}\}$ and variances
$\{\operatorname{Var}_t^{(j)}\}$ to decide whether to continue or stop local
search at this timestep, subject to the budget $K_t$.
The final chosen noise at step $t$ is then fed into the base sampler
$F_\theta$ to obtain $x_{t-1}$, and the process repeats until we reach $x_0$.

Under the same total NFE budget, this two-level scheduling---offline coarse
time allocation plus online fine-grained control based on gain and variance---is
designed to place more computation on timesteps that are both globally and
instance-wise valuable, and to avoid wasting search on steps where noise
perturbations have little impact on the verifier.
In \cref{sec:experiments} we show empirically that this global scheduling
significantly outperforms uniform per-step allocation on ImageNet and
text-to-image benchmarks.

\begin{algorithm}[t]
\caption{Budget Scheduling with Offline Profiling and Online Control}
\label{alg:offline_online_budget}
\begin{algorithmic}[1]
\STATE \textbf{Input:} timesteps $\{t_T,\dots,t_1\}$, total budget $B$, offline allocation $\{K_t\}$ ($\sum_t K_t=B$),
slack $\delta$, gain window $W_g$, threshold coefs $(\beta_g,\beta_\sigma)$, black-box $\textsc{LocalIter}(t)$.
\STATE \textbf{Output:} realized per-step iterations $\{\widehat K_t\}$.

\STATE $R \leftarrow B$; \quad $\mathcal{H}_g\leftarrow\emptyset$ (historical gains); \quad $\mathcal{H}_\sigma\leftarrow\emptyset$ (historical variances).
\FOR{$t=T,T\!-\!1,\dots,1$}
    \STATE $K \leftarrow \min(K_t, R)$; \quad $K_{\max}\leftarrow \min(K+\delta, R)$; \quad $j_{\text{watch}}\leftarrow \max(1, K-\delta)$.
    \STATE $s^{(0)} \leftarrow -\infty$; \quad $\mathcal{H}\leftarrow\emptyset$ (recent gains); \quad $\widehat K_t \leftarrow 0$.
    \FOR{$j=1,2,\dots,K_{\max}$}
        \STATE $(s^{(j)}, \sigma^{(j)}) \leftarrow \textsc{LocalIter}(t)$
        \STATE $g^{(j)} \leftarrow s^{(j)} - s^{(j-1)}$; push $g^{(j)}$ into $\mathcal{H}$ and keep last $W_g$.
        \STATE $\widetilde g^{(j)} \leftarrow \mathrm{Mean}(\mathcal{H})$.
        \STATE $\widehat K_t \leftarrow \widehat K_t + 1$; \quad $R \leftarrow R - 1$.
        \IF{$R=0$} \STATE \textbf{break} \ENDIF
        \IF{$j \ge j_{\text{watch}}$ \AND $|\mathcal{H}_g|>0$ \AND $|\mathcal{H}_\sigma|>0$}
            \STATE $g_{\mathrm{th}}\leftarrow \beta_g \cdot \mathrm{Mean}(\mathcal{H}_g)$;\quad
                   $\sigma_{\mathrm{th}}\leftarrow \beta_\sigma \cdot \mathrm{Mean}(\mathcal{H}_\sigma)$.
            \IF{$\widetilde g^{(j)} < g_{\mathrm{th}}$ \AND $\sigma^{(j)} < \sigma_{\mathrm{th}}$}
                \STATE \textbf{break}
            \ENDIF
        \ENDIF
    \ENDFOR
    \STATE append $\mathrm{Mean}(\mathcal{H})$ to $\mathcal{H}_g$; append $\sigma^{(\widehat K_t)}$ to $\mathcal{H}_\sigma$.
\ENDFOR
\STATE \textbf{return} $\{\widehat K_t\}$.
\end{algorithmic}
\end{algorithm}


\section{Experiments}
\label{sec:experiments}

\subsection{Experimental Setting}
\label{sec:exp-setting}
We evaluate global budget scheduling for verifier-guided noise trajectory search on both Stable Diffusion and EDM under \emph{strict} function-evaluation budgets (NFE).
Our evaluation uses two lightweight verifiers that have been adopted as single-image rewards in prior noise-trajectory-search benchmarks: \textbf{Brightness}, computed from perceived luminance, and \textbf{Compressibility}, computed from the maximally-compressed JPEG file size and normalized to a fixed range.%
\footnote{Because resolution is fixed in our setup, JPEG size is determined by image compressibility.}
These verifiers enable controlled, repeatable measurements of test-time scaling without human annotation.

For the \textbf{baseline}, we follow a standard local search recipe used in prior work: an $\epsilon$-greedy noise search operator combined with a \textbf{uniform per-timestep} compute schedule (``Naive'').
Our method (\textbf{Online (Ours)}) keeps the same family of local search operators but replaces the global policy with a scheduler combining offline profiling and online control: offline profiling produces a coarse per-step allocation, and online control applies windowed early stopping based on (i) observed score gains and (ii) candidate score variance, while enforcing the exact total NFE budget.
All experiments are run on a single NVIDIA A100 GPU; in our implementation, wall-clock generation time is dominated by UNet function evaluations and is therefore effectively determined by the NFE budget.

\subsection{Stable Diffusion: Scaling with NFE}
\label{sec:exp-sd-scaling}
\begin{table}[t]
\centering
\caption{SD results across different NFE budgets.}
\label{tab:sd_budget}
\begin{tabular}{c|c|c}
\toprule
\textbf{NFE} & \textbf{Naive} & \textbf{Online (Ours)} \\
\midrule
\multicolumn{3}{c}{\textit{Brightness}} \\
\midrule
400 & $0.7025 \pm 0.0466$ & $0.7248 \pm 0.0563$ \\
500 & $0.7094 \pm 0.0631$ & $0.7346 \pm 0.0599$ \\
800 & $0.7511 \pm 0.0620$ & $0.7832 \pm 0.0462$ \\
\midrule
\multicolumn{3}{c}{\textit{Compressibility}} \\
\midrule
400 & $0.8833 \pm 0.0166$ & $0.8946 \pm 0.0164$ \\
500 & $0.8825 \pm 0.0170$ & $0.9004 \pm 0.0147$ \\
800 & $0.8892 \pm 0.0198$ & $0.9008 \pm 0.0123$ \\
\bottomrule
\end{tabular}
\end{table}

Table~\ref{tab:sd_budget} shows that our method improves not only absolute
performance but also the \emph{speed} of quality scaling with respect to NFE on
Stable Diffusion.
For \textit{Brightness}, the baseline increases from $0.7025$ (400 NFE) to
$0.7094$ (500 NFE) and $0.7511$ (800 NFE), while our method achieves
$0.7248$, $0.7346$, and $0.7832$ at the same budgets, yielding consistent
margins of +0.0223, +0.0252, and +0.0321, respectively.
Crucially, our 400-NFE result ($0.7248$) already surpasses the baseline at
500 NFE ($0.7094$), meaning we reach a higher brightness level with
\textbf{20\% fewer} function evaluations.
The speed advantage is even clearer for \textit{Compressibility}.
The baseline progresses from $0.8833$ (400) to $0.8825$ (500) and $0.8892$ (800),
whereas our method reaches $0.8946$, $0.9004$, and $0.9008$.
Notably, our method at \textbf{400 NFE} ($0.8946$) already exceeds the baseline
at \textbf{800 NFE} ($0.8892$), achieving better compressibility with
\textbf{50\% fewer} function evaluations.
Similarly, our 500-NFE result ($0.9004$) also surpasses the baseline at 800 NFE,
corresponding to a \textbf{37.5\%} compute reduction for a higher score.
Overall, across all three budgets the quality curve is consistently shifted
leftward, showing that global scheduling with per-step awareness converts
additional NFE into quality more efficiently than uniform allocation.

\subsection{EDM: Scaling with NFE}
\label{sec:exp-edm-scaling}
\begin{table}[t]
\centering
\caption{EDM results across different NFE budgets.}
\label{tab:edm_budget}
\begin{tabular}{c|c|c}
\toprule
\textbf{NFE} & \textbf{Naive} & \textbf{Online (Ours)} \\
\midrule
\multicolumn{3}{c}{\textit{Brightness}} \\
\midrule
144 & $0.9507 \pm 0.0153$ & $0.9887 \pm 0.0046$ \\
180 & $0.9617 \pm 0.0160$ & $0.9904 \pm 0.0036$ \\
288 & $0.9787 \pm 0.0178$ & $0.9988 \pm 0.0012$ \\
\midrule
\multicolumn{3}{c}{\textit{Compressibility}} \\
\midrule
144 & $0.6714 \pm 0.0165$ & $0.6845 \pm 0.0074$ \\
180 & $0.6774 \pm 0.0120$ & $0.7003 \pm 0.0089$ \\
288 & $0.6901 \pm 0.0118$ & $0.7165 \pm 0.0085$ \\
\bottomrule
\end{tabular}
\end{table}
Table~\ref{tab:edm_budget} demonstrates a pronounced acceleration effect on EDM:
our method reaches high quality at substantially lower NFE, indicating a strong
left-shift of the quality--compute curve.
For \textit{Brightness}, the baseline improves from $0.9507$ (144 NFE) to
$0.9617$ (180 NFE) and $0.9787$ (288 NFE), while our method achieves
$0.9887$, $0.9904$, and $0.9988$, respectively (gaps of +0.0380, +0.0287,
and +0.0201).
Importantly, our 144-NFE brightness ($0.9887$) already surpasses the baseline at
288 NFE ($0.9787$), meaning we obtain higher quality with \textbf{50\% fewer}
function evaluations.
Likewise, our 180-NFE result ($0.9904$) also exceeds the baseline at 288 NFE,
corresponding to a \textbf{37.5\%} compute reduction.
These comparisons across the three budgets directly show that our scheduler
achieves the same (or better) quality with substantially fewer denoising
evaluations.
The same acceleration holds for \textit{Compressibility}.
The baseline increases from $0.6714$ (144) to $0.6774$ (180) and $0.6901$ (288),
whereas our method reaches $0.6845$, $0.7003$, and $0.7165$.
Here, our 144-NFE result ($0.6845$) already exceeds the baseline at 180 NFE
($0.6774$), saving \textbf{20\%} NFE for a higher score; more strikingly,
our 180-NFE result ($0.7003$) surpasses the baseline at 288 NFE ($0.6901$),
saving \textbf{37.5\%} NFE.
Across all three budgets, the improvements grow from +0.0131 to +0.0229 and
+0.0264, suggesting that global scheduling continues to allocate search compute
to impactful timesteps as budget increases rather than wasting NFE uniformly.

\subsection{Ablation: Offline Allocation vs.\ Offline+Online Control}
\label{sec:exp-ablation-offline-online}
\begin{table}[t]
\centering
\caption{Stable Diffusion results under fixed NFE $=400$.
\textbf{Naive} uses uniform per-step allocation.
\textbf{Online (Ours)} uses offline profiling plus online control.
\textbf{Offline Only} removes online control and uses only offline allocation.}
\label{tab:sd_400_ablation}
\begin{tabular}{lccc}
\toprule
\textbf{Metric} & \textbf{Naive} & \textbf{Online (Ours)} & \textbf{Offline Only} \\
\midrule
Brightness        & $0.7025 \pm 0.0466$ & $\mathbf{0.7248 \pm 0.0563}$ & $0.7158 \pm 0.0645$ \\
Compressibility   & $0.8833 \pm 0.0166$ & $\mathbf{0.8946 \pm 0.0164}$ & $0.8873 \pm 0.0183$ \\
\bottomrule
\end{tabular}
\end{table}

Table~\ref{tab:sd_400_ablation} isolates the contribution of online control at a fixed compute budget.
The \textbf{Offline Only} variant already improves over uniform allocation on both metrics, showing that coarse per-step budgeting is beneficial by itself.
Adding online feedback and windowed early stopping further improves performance, yielding the best scores for both \textit{Brightness} (0.7248) and \textit{Compressibility} (0.8946).
This indicates that offline profiling captures \emph{where} search tends to matter, while online control refines \emph{how much} of the allocated budget is actually worth spending on a given prompt and timestep.

\subsection{Larger Prompt Set Evaluation}
\label{sec:exp-larger-prompts}
\begin{table}[t]
\centering
\caption{Larger prompt set evaluation.
We use 20 prompts, each repeated 10 times.
Baseline corresponds to the uniform allocation method; Ours is the proposed global scheduler.}
\label{tab:larger_prompt_exp}
\begin{tabular}{lcc}
\toprule
 & \textbf{brightness} & \textbf{compressibility} \\
\midrule
Baseline          & $0.6949 \pm 0.0756$ & $0.7938 \pm 0.0892$ \\
Ours              & $\mathbf{0.7213 \pm 0.0752}$ & $\mathbf{0.8076 \pm 0.0962}$ \\
\bottomrule
\end{tabular}
\end{table}

Table~\ref{tab:larger_prompt_exp} evaluates robustness on a broader set of prompts with repeated trials.
Our scheduler consistently improves the mean scores for both verifiers: brightness increases from 0.6949 to 0.7213, and compressibility increases from 0.7938 to 0.8076.
The improvement persists under repeated sampling, suggesting that the gains are not driven by a small subset of prompts but reflect a systematic advantage of allocating compute to high-impact timesteps and early-stopping low-yield search.

\subsection{Compatibility with Different Local Search Operators}
\label{sec:exp-local-operator}
\begin{table}[t]
\centering
\caption{Global scheduling combined with different local search operators.
Left block uses uniform allocation; right block uses our scheduler.
B: Brightness, C: Compressibility.}
\label{tab:local_search_combo}
\begin{tabular}{lcccc}
\toprule
& \multicolumn{2}{c}{\textbf{Naive}} & \multicolumn{2}{c}{\textbf{Ours}} \\
\cmidrule(lr){2-3} \cmidrule(lr){4-5}
\textbf{Metric} & \textbf{Zero-order} & \textbf{Random} & \textbf{Zero-order} & \textbf{Random} \\
\midrule
\textbf{B} & $0.4920 \pm 0.0723$ & $0.7382 \pm 0.0392$ & $\mathbf{0.5080 \pm 0.0752}$ & $\mathbf{0.7457 \pm 0.0452}$ \\
\textbf{C} & $0.5727 \pm 0.0660$ & $0.8417 \pm 0.0268$ & $\mathbf{0.5832 \pm 0.0613}$ & $\mathbf{0.8454 \pm 0.0252}$ \\
\bottomrule
\end{tabular}
\end{table}

Table~\ref{tab:local_search_combo} shows that global scheduling is complementary to the choice of local search operator.
Under both \textbf{Zero-order} and \textbf{Random} local search, replacing a uniform schedule with our scheduler improves \textit{Brightness} and \textit{Compressibility}.
This supports the intended modularity of our two-level formulation: the scheduler targets \emph{where/when} to spend compute, and can be paired with different per-timestep proposal mechanisms without changing the overall framework.

\subsection{Applicability to Flow-based Models}
\label{sec:exp-flow}

To demonstrate that our scheduling framework applies beyond stochastic samplers,
we evaluate on a flow-based model that uses deterministic ODE transitions.
Following prior work~\citep{black2023training}, we introduce stochasticity by
injecting small noise at selected timesteps, enabling exploration of different
generation paths that would otherwise be inaccessible with purely deterministic
sampling.

\begin{table}[t]
\centering
\caption{Flow-based model results (DDIM with $\eta=0$, noise injected for exploration).
Even with originally deterministic sampling, our global scheduling improves quality
by strategically allocating search compute.}
\label{tab:flow_results}
\begin{tabular}{c|c|c}
\toprule
\textbf{NFE} & \textbf{Naive} & \textbf{Online (Ours)} \\
\midrule
\multicolumn{3}{c}{\textit{Brightness}} \\
\midrule
200 & $0.6842 \pm 0.0512$ & $0.7089 \pm 0.0478$ \\
400 & $0.7156 \pm 0.0489$ & $0.7401 \pm 0.0432$ \\
\midrule
\multicolumn{3}{c}{\textit{Compressibility}} \\
\midrule
200 & $0.8621 \pm 0.0198$ & $0.8789 \pm 0.0175$ \\
400 & $0.8803 \pm 0.0187$ & $0.8954 \pm 0.0162$ \\
\bottomrule
\end{tabular}
\end{table}

Table~\ref{tab:flow_results} shows that global scheduling provides consistent
improvements even when applied to deterministic ODE-based sampling with injected
stochasticity for exploration.
At 200 NFE, our method improves \textit{Brightness} by +0.0247 and
\textit{Compressibility} by +0.0168.
At 400 NFE, the improvements are +0.0245 and +0.0151 respectively.
This validates our claim from the introduction: the deliberate introduction of
noise for exploration enables our scheduling framework to work with any sampler
architecture, not just inherently stochastic ones.
The consistent gains across both NFE budgets demonstrate that our scheduler
effectively identifies high-value timesteps even when the underlying sampler
is deterministic.


\section{Conclusion}
\label{sec:conclusion}

We study global scheduling with per-step awareness for inference-time noise trajectory search in diffusion models under strict NFE budgets.
By separating local refinement from global allocation, we design a scheduler combining offline profiling and online control that uses windowed gain and variance signals for early stopping while enforcing an exact budget.
Across Stable Diffusion and EDM, the scheduler consistently improves scores and accelerates quality scaling, often matching higher-NFE baselines with substantially fewer evaluations.

\bibliographystyle{plainnat}
\bibliography{main}


\end{document}
